% Part: computability
% Chapter: computability-theory
% Section: ce-sets

\documentclass[../../../include/open-logic-section]{subfiles}

\begin{document}

\olfileid{cmp}{thy}{ces}
\olsection{संगणनक्षमपणे प्रगणनीय संच}

\begin{defn}
एखादा संच रिकामा असेल किंवा एखाद्या संगणनक्षम फलनाची व्याप्ती असेल,
तर तो \emph{संगणनक्षमपणे प्रगणनीय} असतो.
\end{defn}

\begin{history}
संगणनक्षमपणे प्रगणनीय संचांना \emph{पुनरावर्तीपणे प्रगणनीय}
असेही म्हणतात. ही मूळ संज्ञा आहे. आज ही दोन्ही नावे, तसेच
``c.e.'' आणि ``r.e.'' ही संक्षिप्त रूपे, सर्वसाधारणपणे वापरली जातात.
\end{history}

\begin{explain}
या व्याख्येचा अर्थ काय आणि ही संज्ञा योग्य का आहे, याचा विचार करा.
कल्पना अशी की $S$ हा संगणनक्षम फलन~$f$ ची व्याप्ती असेल, तर
\[
S = \{ f(0), f(1), f(2), \dots \},
\]
म्हणून $f$ हे~$S$ च्या घटकांचे ``प्रगणन'' करते असे मानता येते.
व्याख्येनुसार $f$ वाढते फलन असणे आवश्यक नाही, म्हणजे प्रगणन
चढत्या क्रमाने असणे आवश्यक नाही. किंबहुना, $f$ एकास-एक असणेही
आवश्यक नाही; म्हणजे $f(0)$, $f(1)$, $f(2)$, \dots{} या~$S$
च्या प्रगणनात पुनरुक्ती चालते. उदाहरणार्थ, स्थिर फलन $f(x) = 0$
हे संच $\{ 0
\}$ चे प्रगणन करते.
\end{explain}

प्रत्येक संगणनक्षम संच संगणनक्षमपणे प्रगणनीय असतो. हे पाहण्यासाठी,
$S$~संगणनक्षम आहे असे समजा. $S$ रिकामा असेल, तर व्याख्येवरून
तो संगणनक्षमपणे प्रगणनीय आहे. अन्यथा $a$ हा $S$ मधील कोणताही
घटक घ्या. $f$ ची पुढीलप्रमाणे व्याख्या करा:
\[
f(x) =
\begin{cases}
x & \text{जर $\Char{S}(x) = 1$ असेल तर} \\
a & \text{अन्यथा.}
\end{cases}
\]
मग $f$ हे संगणनक्षम फलन आहे आणि $S$ ही~$f$ ची व्याप्ती आहे.

\end{document}
